\section{Discussion and limitations}
\label{sec:discussion}
\paragraph{What the evidence supports.} With labels, unlabeled data and predictor fixed, the room for shorter valid CLT intervals for a mean is the nonlinearity of $E[Y\mid f]$ beyond its best affine fit, plus whatever PPI++'s clip of $\hat\lambda$ discards. On the real datasets studied the first is below $1\%$ of the width, the second matters only for gene expression and census income, and the estimation cost of a calibration map outweighs it at every $n\le200$; only the synthetic nonlinear populations show a benefit, and it stayed above the ceiling in all 45 cells. We therefore read the outcome as \emph{conditional}: possible gains exist but need a predictor whose conditional mean is clearly non-affine in its own output, e.g.\ a systematically miscalibrated or non-monotone score, and enough labels to learn the map (hundreds). Absent that, PPI++ is already near the ceiling.

\paragraph{Limitations.} (a)~\emph{Estimand and information set:} only the population mean (prevalence for binary $Y$); quantiles and regression coefficients are not studied. Covariates $X$ are not used; the ceiling would be lower if they were, and we did not measure that headroom on the census sets. (b)~\emph{Sampling model:} the pool is the population and records are drawn with replacement, so finite-population effects of the original data protocols are not studied; a without-replacement forest ablation shows no clear difference at $n=50$ and $100$ and higher coverage at $n=200$ ($0.952$ against $0.930$), as the finite-population correction would give (Appendix Table~\ref{tab:wor}). Cells at different $n$ share one pool and are not independent. (c)~\emph{Methods:} one calibrator class (isotonic, $K=5$) was frozen before the final runs; it cannot fit non-monotone maps, and the spline ablation shows a larger gain in that case, but spline was not pre-selected. Learning $\hat g$ also reduces effective sample size at small $n$; we do not claim the best possible estimator of $g$. (d)~\emph{Evidence:} the Monte-Carlo standard error of a coverage is about $0.005$ ($95\%$ half-width about $0.01$); we report cells below nominal rather than testing them and do not correct for multiplicity, so single-cell shortfalls should not be over-read, only the consistent patterns across $n$ and datasets. The cause of census-income undercoverage is untested. (e)~\emph{Theory:} Proposition~\ref{prop:floor} is asymptotic; it does not describe finite-$n$ coverage or the plug-in cost, which we only measure. (f)~\emph{Design:} most cells use $N=20n$; the $N/n$ sweep covers four real sets at $n=100,400$. The conformal baseline is not simulated (it is infinite here). Ballots have no labeled unlabeled split, so its pool is the labeled set. (g)~\emph{Shift:} the plankton result is one real shift; we explain it through $\hat\lambda$ but did not test remedies.

\paragraph{Broader impact.} Intervals that undercover silently mislead downstream science, which is the risk behind the small-$n$ and shift findings; we release only analysis code on public data.
